% !TeX program = pdflatex
% The Pentagon, Machine-Checked — Abel's five-term relation for the dilogarithm in Lean 4.
% Build: pdflatex five-term-lean.tex (twice, for refs)
% Verified against the Lean source 2026-06-13 (lake build Dilog.FiveTerm, exit 0;
% #print axioms on every named theorem -> propext, Classical.choice, Quot.sound only).
% Toolchain v4.30.0-rc2, Mathlib pin 701fb6e9.

\documentclass[11pt]{article}
\usepackage[T1]{fontenc}
\usepackage{lmodern}
\DeclareUnicodeCharacter{2082}{\ensuremath{_2}}
\usepackage{amsmath,amssymb,amsthm,booktabs,graphicx,hyperref}
\usepackage{xurl}
\usepackage[table]{xcolor}
\definecolor{ggteal}{HTML}{1B6F8C}
\definecolor{ggband}{HTML}{DCE7EB}
\definecolor{ggamber}{HTML}{E08A1E}
\usepackage[margin=1in]{geometry}
\usepackage{microtype}
\usepackage{float}
\usepackage[section]{placeins}
\usepackage[most]{tcolorbox}
\usepackage{tikz}
\usetikzlibrary{arrows.meta,positioning}
\setlength{\emergencystretch}{2em}
\newtheorem{theorem}{Theorem}
\newtheorem{lemma}{Lemma}
\newtheorem{remark}{Remark}
\newcommand{\lean}[1]{\texttt{#1}}
\newcommand{\R}{\mathbb{R}}
\newcommand{\Li}{\mathrm{Li}_2}
\DeclareMathOperator{\Lrog}{L}
\hypersetup{colorlinks=true,linkcolor=ggteal,citecolor=ggteal,urlcolor=ggteal}
\newtcbox{\keybox}{on line,boxrule=0.9pt,colframe=ggteal,colback=ggband!40,
  arc=3pt,left=8pt,right=8pt,top=5pt,bottom=5pt}
\newcommand{\keyeq}[1]{\begin{center}\keybox{$\displaystyle #1$}\end{center}}

\title{\bf The Pentagon, Machine-Checked\\[4pt]
  \large\normalfont Abel's five-term relation for the dilogarithm in Lean~4}
\author{Carles Mar\'in\thanks{Formalized with AI assistance (Claude, Anthropic); the
mathematics and all claims are the author's responsibility. The exact boundary of the
contribution is stated in the abstract and Section~\ref{sec:boundary}.}\\[3pt]
  \normalsize Independent researcher \quad\textbar\quad \texttt{karlesmarin@gmail.com}\\
  \normalsize ORCID 0009-0007-5637-9688}
\date{June 2026}

\begin{document}
\maketitle

\begin{abstract}
I give a machine-checked, \lean{sorry}-free proof, in Lean~4 over Mathlib, of
\emph{Abel's five-term relation} for the dilogarithm. With $\Lrog$ the Rogers
$\Lrog$-function, it reads, for $x,y\in(0,1)$,
\[
  \Lrog(x)+\Lrog(y)=\Lrog(xy)+\Lrog\!\Big(\tfrac{x(1-y)}{1-xy}\Big)
    +\Lrog\!\Big(\tfrac{y(1-x)}{1-xy}\Big).
\]
This single identity is the whole weight-2 theory in one line: by theorems of Wojtkowiak
and de~Jeu, \emph{every} functional equation of $\Li$ with rational arguments is a
consequence of it, and it is the defining relation of the Bloch group and the pentagon of
cluster algebras. The proof is the classical differentiation argument made fully rigorous:
fixing $y$, the five-term difference has identically zero derivative on $(0,1)$ -- nine
logarithmic terms collapse into a basis of five and the rest is \lean{ring} -- so it is
constant, and its boundary value as $x\to0^+$ is computed by squeezing $\log
x\cdot\log(1-x)\to0$. The development is self-contained: Mathlib carries no dilogarithm, so
this sits on a from-scratch Lean construction of $\Li$. I am exact about the boundary: the
mathematics is entirely classical (Spence, Abel, Rogers; Wojtkowiak's and de~Jeu's
universality); the contribution is the formalization -- to the best of my knowledge the
first machine-checked proof of the five-term relation, and the first machine-checked
dilogarithm, in any proof assistant -- together with three reusable building blocks. The
headline theorem and its blocks depend only on \lean{propext}, \lean{Classical.choice} and
\lean{Quot.sound}.
\end{abstract}

\section{One identity to rule them all}\label{sec:intro}

Put $x=y=\tfrac12$. The five-term relation collapses, after $\tfrac{(1/2)(1/2)}{1-1/4}
=\tfrac13$, to
\[
  2\,\Lrog\!\big(\tfrac12\big)=\Lrog\!\big(\tfrac14\big)+2\,\Lrog\!\big(\tfrac13\big),
  \qquad\text{i.e.}\qquad \Lrog\!\big(\tfrac14\big)+2\,\Lrog\!\big(\tfrac13\big)=\frac{\pi^2}{6},
\]
since $\Lrog(\tfrac12)=\pi^2/12$. That a sum of three dilogarithm values at $\tfrac14$ and
$\tfrac13$ should be exactly $\pi^2/6$ is not obvious; it is one of infinitely many
shadows of a single two-variable law. The dilogarithm $\Li(x)=\sum_{n\ge1}x^n/n^2$ obeys
a handful of functional equations. Two are elementary -- Euler's \emph{reflection}
$\Li(x)+\Li(1-x)=\tfrac{\pi^2}{6}-\log x\log(1-x)$ and \emph{Landen}'s transformation --
and one is not: the \emph{five-term relation}, found by Spence (1809) and Abel (1830) and
rediscovered ever since. It is the deepest weight-2 identity, and the only one that truly
needs two variables.

Its centrality is a theorem, not a slogan. Wojtkowiak proved that every one-variable
functional equation of $\Li$ with rational arguments is a $\mathbb{Z}$-linear combination
of instances of the five-term relation; de~Jeu extended this to any number of
variables~\cite{Wojtkowiak,deJeu}. The same identity \emph{defines} the Bloch group, so it
governs scissors-congruence invariants and the volumes of hyperbolic $3$-manifolds
\cite{Zagier}; in cluster algebras it is the \emph{pentagon}, the relation of five
mutations~\cite{Nakanishi}. Write the five arguments on a cycle
(Figure~\ref{fig:pentagon}) and the five-fold symmetry that powers all of this appears.

The cleanest statement uses the \emph{Rogers $\Lrog$-function}
$\Lrog(x):=\Li(x)+\tfrac12\log x\log(1-x)$, which absorbs the stray logarithms of the raw
$\Li$ form and turns the relation into a pure sum of five $\Lrog$ values:
\keyeq{\Lrog(x)+\Lrog(y)=\Lrog(xy)+\Lrog\!\Big(\tfrac{x(1-y)}{1-xy}\Big)
  +\Lrog\!\Big(\tfrac{y(1-x)}{1-xy}\Big),\qquad x,y\in(0,1).}

\begin{figure}[t]
\centering
\begin{tikzpicture}[scale=1.5,every node/.style={font=\small}]
  \foreach \i/\lab in {0/{$x$},1/{$\tfrac{x(1-y)}{1-xy}$},2/{$y$},
      3/{$\tfrac{y(1-x)}{1-xy}$},4/{$xy$}}{
    \node (n\i) at ({90+\i*72}:1.25) {};
    \node at ({90+\i*72}:1.62) {\lab};
    \fill[ggteal] ({90+\i*72}:1.25) circle (1.5pt);
  }
  \foreach \i in {0,...,4}{
    \pgfmathtruncatemacro{\j}{mod(\i+1,5)}
    \draw[ggteal,-{Stealth[length=4pt]}] (n\i) -- (n\j);
  }
\end{tikzpicture}
\caption{\label{fig:pentagon}The five arguments on a cycle. The five-fold symmetry of this
pentagon is the combinatorial shadow of the Bloch-group relation and of the
cluster-algebra mutation pentagon.}
\end{figure}

\paragraph{The state of formalization.} Mathlib, the largest formal library, carries no
dilogarithm: neither $\Li$ nor any of its functional equations. A check of the Lean,
Isabelle and Rocq (Coq) libraries in June~2026 found no formalization of the five-term
relation, nor of $\Li$ as an analytic object beyond elementary special values. What
follows therefore sits on a from-scratch Lean construction of $\Li$ -- its series,
derivative $-\log(1-x)/x$, reflection, Landen, duplication -- of which the five-term
relation is the capstone. With it, the two-variable identities and the polylogarithm
\emph{ladders} (the golden-ratio and quartic-base evaluations) become corollaries rather
than separate formalization targets.

\section{The statement, and the idea of the proof}\label{sec:proof}

\begin{theorem}[\lean{Dilog.rogersL\_fiveterm}]\label{thm:five}
For all $x,y\in(0,1)$,
\[
  \Lrog(x)+\Lrog(y)=\Lrog(xy)+\Lrog\!\Big(\tfrac{x(1-y)}{1-xy}\Big)
    +\Lrog\!\Big(\tfrac{y(1-x)}{1-xy}\Big).
\]
\end{theorem}

I prove it the textbook way -- by differentiation -- but with every step made rigorous.
Fix $y\in(0,1)$ and set
\[
  G(t)=\Lrog(t)-\Lrog(ty)-\Lrog\!\Big(\tfrac{t(1-y)}{1-ty}\Big)
    -\Lrog\!\Big(\tfrac{y(1-t)}{1-ty}\Big).
\]
The theorem is $G(x)=-\Lrog(y)$. Three movements.

\paragraph{1. The derivative vanishes.} For $x,y\in(0,1)$ all three composite arguments
again lie in $(0,1)$, so every term is differentiable. With $\Lrog'(t)=-\tfrac12\big(
\tfrac{\log(1-t)}{t}+\tfrac{\log t}{1-t}\big)$, the chain rule writes $G'(t)$ as four such
expressions weighted by the derivatives of the composite arguments. Every logarithm that
occurs -- of $t$, of $ty$, of the two fractional arguments and their complements --
expands, by $\log(ab)=\log a+\log b$ and $\log(a/b)=\log a-\log b$, into the five-element
basis $\{\log t,\log y,\log(1-t),\log(1-y),\log(1-ty)\}$. In that basis $G'(t)$ is a
rational expression whose every coefficient cancels: $G'\equiv0$. In Lean this is the one
delicate step (\lean{fiveterm\_hasDerivAt\_zero}): the nine logarithms are rewritten into
the basis and the rational remainder is cleared and closed by \lean{ring}. (The
cancellation was checked numerically to thirty digits before it was formalized -- and the
formalization caught a subtlety the eye misses: \lean{field\_simp} silently reorders
$x\cdot y$ to $y\cdot x$ and then needs $1-y x\neq0$ in hand to discharge the resulting
inverse, a fact that simply had to be in scope at the right moment.)

\paragraph{2. Hence $G$ is constant.} A function with vanishing derivative on the open
connected interval $(0,1)$ is constant there -- Mathlib's
\lean{IsOpen.is\_const\_of\_deriv\_eq\_zero}, the mean-value theorem in disguise. So
$G(x)=G(t)$ for every $t\in(0,1)$.

\paragraph{3. The boundary value.} As $t\to0^+$, $ty\to0$, $\tfrac{t(1-y)}{1-ty}\to0$ and
$\tfrac{y(1-t)}{1-ty}\to y$, so $G(t)\to\Lrog(0)-\Lrog(0)-\Lrog(0)-\Lrog(y)=-\Lrog(y)$ --
\emph{provided} $\Lrog$ is right-continuous at $0$ with $\Lrog(0)=0$. This is the one
analytic subtlety. $\Lrog(t)=\Li(t)+\tfrac12\log t\log(1-t)$, and while $\Li$ is
continuous, $\log t\cdot\log(1-t)$ is an indeterminate $(-\infty)\cdot 0$. It is tamed by
a squeeze: on $(0,\tfrac12]$ one has $-\log(1-t)\le 2t$ (from $\log z\ge1-z^{-1}$), so
\[
  0\ \le\ \log t\cdot\log(1-t)=(-\log t)\,(-\log(1-t))\ \le\ 2t\,(-\log t)=2\,(-t\log t)
  \ \longrightarrow\ 0,
\]
because $t\log t\to0$ -- Mathlib's \lean{Real.continuous\_mul\_log}. Thus $G\to-\Lrog(y)$;
being constant and equal to $G(x)$ near $0$, it \emph{equals} $-\Lrog(y)$, which is
Theorem~\ref{thm:five}. \qed

\begin{tcolorbox}[colframe=ggamber,colback=ggamber!8,arc=3pt,boxrule=0.8pt,
  title=\textbf{The recipe: read off a $\pi^2$-evaluation in three lines},
  coltitle=white,colbacktitle=ggamber]
Pick any $x,y\in(0,1)$ at which the four arguments are nice. (1) Form
$xy,\ \tfrac{x(1-y)}{1-xy},\ \tfrac{y(1-x)}{1-xy}$. (2) Apply
$\Lrog(x)+\Lrog(y)=$ their sum, with a machine-checked guarantee. (3) Substitute any
$\Lrog$ values you know.\\[3pt]
\emph{Worked:} take $x=y=\tfrac12$. Then $xy=\tfrac14$ and both fractions are $\tfrac13$,
so $2\,\Lrog(\tfrac12)=\Lrog(\tfrac14)+2\,\Lrog(\tfrac13)$. Since $\Lrog(\tfrac12)=
\Li(\tfrac12)+\tfrac12\log^2\!2$ and $\Li(\tfrac12)=\tfrac{\pi^2}{12}-\tfrac12\log^2\!2$,
the log terms cancel and $\Lrog(\tfrac12)=\tfrac{\pi^2}{12}$; hence
$\Lrog(\tfrac14)+2\,\Lrog(\tfrac13)=\tfrac{\pi^2}{6}$ -- a closed form for three
dilogarithm values, certified.
\end{tcolorbox}

\section{The building blocks}\label{sec:blocks}

The headline rests on three reusable pieces, each \lean{sorry}-free with the standard
three axioms.

\begin{lemma}[\lean{Dilog.hasDerivAt\_rogersL}]
On $(0,1)$, $\displaystyle\Lrog'(x)=-\tfrac12\Big(\tfrac{\log(1-x)}{x}+\tfrac{\log
x}{1-x}\Big)$.
\end{lemma}
The same derivative scheme already used for Euler reflection and Landen: combine
$\Li'(x)=-\log(1-x)/x$ with the product rule for $\tfrac12\log x\log(1-x)$.

\begin{lemma}[\lean{Dilog.tendsto\_rogersL\_nhdsGT\_zero}]
$\Lrog(t)\to0$ as $t\to0^+$.
\end{lemma}
The $\Li$ part is continuous on $[-1,1]$; the log-product is squeezed to $0$ by
$2(-t\log t)$. This is the only place the proof leaves pure algebra for analysis.

\begin{lemma}[\lean{Dilog.continuousAt\_rogersL}]
$\Lrog$ is continuous at every $x\in(0,1)$.
\end{lemma}
Needed for the one surviving boundary term, $\Lrog\big(\tfrac{y(1-t)}{1-ty}\big)\to
\Lrog(y)$.

\section{The boundary of the contribution}\label{sec:boundary}

Everything here that is mathematics is classical. The five-term relation is Spence's and
Abel's; the Rogers $\Lrog$-function is Rogers's (1907); the universality statements are
Wojtkowiak's and de~Jeu's; the differentiation proof is standard. I claim no new identity
and no new inequality. The contributions are three: the \emph{formalization} (to my
knowledge the first machine-checked five-term relation, on the first machine-checked
dilogarithm, in any proof assistant); the \emph{reusable infrastructure} -- the
Rogers-$\Lrog$ derivative, the right-continuity at $0$ via the $-t\log t$ squeeze, and the
interior continuity -- which is exactly what any further weight-2 development (ladders, the
Bloch group, the Bloch--Wigner five-term) will need; and the \emph{proof-engineering
record}, where the nine-into-five logarithmic cancellation is the kind of step a formal
check earns its keep on.

\emph{On methodology.} The formalization was done with an AI assistant (Claude, Anthropic),
with the central cancellation verified numerically to thirty digits before formalizing and
every named theorem checked by \lean{\#print axioms} to depend only on \lean{propext},
\lean{Classical.choice} and \lean{Quot.sound}. Correctness does not rest on the assistant:
the Lean kernel checks every proof. The mathematics and the claims are the author's
responsibility.

\section{Open roads}\label{sec:roads}

\begin{enumerate}
\item \textbf{Ladders as corollaries.} With the five-term relation, the golden-ratio
ladders $\Lrog(1/\varphi)=\pi^2/10$, $\Lrog(1/\varphi^2)=\pi^2/15$ and their quartic-base
relatives are linear-algebra consequences rather than separate proofs; mechanizing this in
Lean is the natural next step.
\item \textbf{The Bloch--Wigner five-term.} The same relation holds for the single-valued
Bloch--Wigner function on $\mathbb{C}$, where it underlies the hyperbolic-volume and
scissors-congruence applications; formalizing it is the complex-analytic sequel.
\item \textbf{Universality.} Wojtkowiak's reduction of every rational one-variable
equation to the five-term would be a deep formalization target, of which this is the
indispensable first brick.
\end{enumerate}

\section*{Reproducing}
\begin{verbatim}
git clone https://github.com/karlesmarin/godsil-gutman-lean
cd godsil-gutman-lean && lake exe cache get
lake build Dilog.FiveTerm
\end{verbatim}
Lean \lean{v4.30.0-rc2}, Mathlib pin \lean{701fb6e9}. The relation is
\lean{Dilog.rogersL\_fiveterm} in \lean{Dilog/FiveTerm.lean}; its supports are
\lean{hasDerivAt\_rogersL}, \lean{fiveterm\_hasDerivAt\_zero},
\lean{tendsto\_rogersL\_nhdsGT\_zero} and \lean{continuousAt\_rogersL}. Each is checked by
\lean{\#print axioms} to depend only on \lean{propext}, \lean{Classical.choice},
\lean{Quot.sound}.

\begin{thebibliography}{9}
\bibitem{Wojtkowiak} Z.\ Wojtkowiak, \emph{The basic structure of polylogarithmic
functional equations}, in \textit{Structural Properties of Polylogarithms} (L.\ Lewin,
ed.), Math.\ Surveys Monogr.\ \textbf{37}, AMS (1991), 205--231.
\bibitem{deJeu} R.\ de Jeu, \emph{Towards regulator formulae for the $K$-theory of curves
and number fields}, Compositio Math.\ \textbf{124} (2000), 137--194.
\bibitem{Zagier} D.\ Zagier, \emph{The dilogarithm function}, in \textit{Frontiers in
Number Theory, Physics, and Geometry II}, Springer (2007), 3--65.
\bibitem{Nakanishi} T.\ Nakanishi, \emph{Dilogarithm identities for conformal field
theories and cluster algebras}, in \textit{Symmetries, Integrable Systems and
Representations}, Springer Proc.\ Math.\ Stat.\ \textbf{40} (2013), 407--444.
\bibitem{Rogers} L.\ J.\ Rogers, \emph{On function sum theorems connected with the series
$\sum x^n/n^2$}, Proc.\ London Math.\ Soc.\ \textbf{4} (1907), 169--189.
\bibitem{Lewin} L.\ Lewin, \emph{Polylogarithms and Associated Functions},
North-Holland (1981).
\bibitem{Mathlib} The Mathlib Community, \emph{The Lean Mathematical Library},
CPP 2020, 367--381.
\end{thebibliography}

\end{document}
